\pdfvariable suppressoptionalinfo 611
\providecommand{\CRevisionRound}{2}
\documentclass[9pt]{extarticle}
\usepackage[a4paper,margin=0.65in]{geometry}
\usepackage{fontspec}
\setmainfont{Latin Modern Roman}
\setsansfont{Latin Modern Sans}
\setmonofont{Latin Modern Mono}
\usepackage[english]{babel}
\babelprovide[import]{chinese-simplified}
\babelfont[chinese-simplified]{rm}[Renderer=Harfbuzz,Language=Default]{Droid Sans Fallback}
\usepackage{amsmath,amssymb,booktabs}
\newcommand{\Rmax}{\mathbb{Q}_{\max}}
\title{Irreducible Max-Plus Powers: Exact Cyclicity, CSR Phases,\\
Orbit Strata, and the Sharp Transient Boundary}
\author{Route-A structural certificate C188}
\date{}
\begin{document}
\maketitle
\begin{abstract}
For every irreducible rational max-plus matrix $A$, normalize by its maximum
cycle mean and let $\gamma$ be the lcm of the cyclicities of the critical
components.  The source-locked cyclicity theorem says that $\gamma$ is the
minimal ultimate period of the entire matrix-power sequence.  We integrate
this with the exact first-equality transient, ultimate $CS^tR$ representation,
raw and projective orbit strata, attraction cones, the eigencone, and periodic
ultimate spans.  A fixed-support $2\times2$ family has every positive transient,
while reducible matrices may have several growth rates.  The classification is
exact but has no intrinsic rational-prime semantics, target divisor, or
Hilbert-space quantization.
\end{abstract}
\noindent\textbf{Keywords:} max-plus algebra; cyclicity; critical graph; CSR;
attraction cone; transient.

\begin{otherlanguage}{chinese-simplified}
\begin{center}{\large 中文摘要}\end{center}
本文对任意不可约有理\foreignlanguage{english}{max-plus}矩阵给出统一的最终动力学分类。
减去最大循环均值后，临界图各强连通分量循环性的最小公倍数$\gamma$，正是整个矩阵幂
序列的最小最终周期。本文将这一有明确来源的经典定理与首次等式暂态、
\foreignlanguage{english}{CSR}表示、向量及射影轨道分层、吸引锥、特征锥和最终列张成
空间联系起来。固定的二阶支撑可以实现任意大的暂态；可约系统则可能具有多个增长率。
这些精确热带结构不产生内生的有理素数语义、目标除子或希尔伯特空间量子化。

\noindent 关键词：\foreignlanguage{english}{max-plus}代数；循环性；临界图；
\foreignlanguage{english}{CSR}；吸引锥；暂态。
\end{otherlanguage}

\section{Frozen system and classical input}
On $\Rmax=\mathbb Q\cup\{-\infty\}$ write
$u\oplus v=\max(u,v)$ and $u\otimes v=u+v$.  Let
$A\in\Rmax^{n\times n}$ have strongly connected finite support.  Its maximum
cycle mean is $\lambda(A)$, and $B=A-\lambda(A)$ means that $\lambda(A)$ is
subtracted from every finite entry.  The cited sources primarily use the
equivalent max-times convention; the logarithm preserves powers, support,
critical cycles and cyclicity while converting it exactly to this max-plus
notation.  The critical graph is the union of all
cycles attaining $\lambda(A)$.  If its strongly connected components have
cyclicities $\gamma_1,\ldots,\gamma_s$---the gcds of their cycle lengths---put
\[
 \gamma=\operatorname{lcm}(\gamma_1,\ldots,\gamma_s).
\]
The classical cyclicity theorem, in the convention source-locked to
Sergeev~\cite{Sergeev09}, states more than a period bound: the \emph{least}
ultimate period of $\{B^t\}_{t\ge0}$ is exactly $\gamma$.  Thus there is a
finite $T_0$ such that
\begin{equation}\label{eq:period}
 B^{t+\gamma}=B^t\qquad(t\ge T_0),
\end{equation}
and no smaller positive integer is an ultimate matrix period.  When the
critical graph has several components the lcm, not a gcd across components,
is essential.

The second source lock is the ultimate CSR theory of Sergeev and
Schneider~\cite{SS12}.  Put
\[
 K=(B^\gamma)^*=I\oplus B^\gamma\oplus\cdots\oplus B^{(n-1)\gamma}.
\]
Let $C$ retain the critical columns of $K$, $R$ its critical rows, and $S$ the
critical edges of $B$, filling all other entries by $-\infty$.  For a
matrix-dependent finite $T_{\rm CSR}$,
\begin{equation}\label{eq:csr}
 B^t=C\otimes S^t\otimes R\qquad(t\ge T_{\rm CSR}).
\end{equation}
Neither source asserts a weight-independent uniform transient, and neither do
we.

\section{Complete eventual dynamics}
\paragraph{Exact transient.}
Define $T(A)$ to be the least onset in~\eqref{eq:period}.  It has the sharp
one-equality characterization
\begin{equation}\label{eq:transient}
 T(A)=\min\{t\ge0:B^{t+\gamma}=B^t\}.
\end{equation}
Indeed, the classical theorem makes the set nonempty.  If equality holds at
one $t$, right multiplication by $B^k$ gives
$B^{t+k+\gamma}=B^{t+k}$ for every $k\ge0$, proving~\eqref{eq:transient}.

\paragraph{Vectors, projective classes, and attraction cones.}
For every $x\in\Rmax^n$, equation~\eqref{eq:period} shows that the ultimate
period of $B^t\otimes x$ divides $\gamma$.  For a nonzero vector, max-plus
homogeneity gives the same conclusion after projectivization.  Both periods
may be smaller: for
\[
 B=\begin{pmatrix}-\infty&0\\0&-\infty\end{pmatrix},
 \qquad x=(0,0)^T,
\]
the matrix period is two whereas $x$ is fixed.

For $p\mid\gamma$, the exact attraction cone is
\[
 \operatorname{Attr}_p(B)
 =\{x:B^{T+p}\otimes x=B^T\otimes x\}.
\]
The same propagation proof shows that it consists precisely of vectors whose
ultimate period divides $p$.  Hence the exact-period-$p$ stratum is
\[
 \operatorname{Attr}_p(B)\setminus
 \bigcup_{\substack{q\mid p\\q<p}}\operatorname{Attr}_q(B).
\]
Replacing equality by projective equality gives the projective strata.  These
are exact two-sided max-plus linear tests, not labels inferred from sampled
orbits.

\paragraph{Eigencone and ultimate spans.}
The normalized eigencone $E(B)=\{x:B\otimes x=x\}$ lies in
$\operatorname{Attr}_1(B)$.  For $0\le r<\gamma$, put
$V_r=\operatorname{Col}(B^{T+r})$.  Then
\[
 V_{r+\gamma}=V_r,\qquad B(V_r)=V_{r+1},
\]
and every orbit lies in the corresponding phase cone after time $T$.  We do
not assert that all $V_r$ are distinct.

\section{Sharp boundaries}
\paragraph{Primitive does not mean immediate.}
If $\gamma=1$, the matrix powers are eventually constant.  Nevertheless, for
every integer $m\ge1$ the fixed-support irreducible family
\[
 B_m=\begin{pmatrix}0&-m\\0&-1\end{pmatrix}
\]
has critical cyclicity one and, for $t\ge1$, direct induction gives
\begin{equation}\label{eq:family}
 B_m^t=\begin{pmatrix}0&-m\\0&\max(-t,-m)\end{pmatrix}.
\end{equation}
Thus its minimal transient is exactly $m$.  Fixed dimension and fixed support
therefore admit unbounded transient as the weights vary; a dimension-only,
weight-independent bound is false.

\paragraph{Reducibility changes the theorem.}
The reducible matrix $D=\operatorname{diag}(0,1)$ has component growth rates
zero and one.  After normalization by the larger rate, its powers contain
$-t$ and are not a single periodic sequence.  General reducible behavior is
described by multiple growth rates and CSR terms~\cite{SS12}; we do not extend
the irreducible conclusion across this boundary.

\section{Exact evidence and Route-A stop}
The exact ledger covers 177 matrices, 901 raw/projective vector rows, 441
simple cycles and 189 critical components.  It checks 5,469 CSR cells and
7,471 propagated-period cells.  An independent Karp/closure/Tarjan/gcd and
binary-power checker passes 7,924 assertions.  A separate SymPy path passes
10,615 checks; replay is byte exact; 137 repaired-hash and one stale-hash
attacks are rejected.  Among the regression vectors, 28 raw and 28 projective
orbits have period strictly below their matrix period.

These checks are regression, not a finite proof of the source theorems.  The
package contribution is the source-locked synthesis, elementary derived
classification and executable boundary certificate.  Rational weights,
support cycles and critical means carry no intrinsic rational-prime origin;
ultimate phases and CSR identify no target divisor; no target functional
equation, continuation, counting law or Weil compression follows; and a
max-plus semimodule is not a source-native complex Hilbert-space quantization.
Therefore
\[
 (A0,A1,A2,A3,A4)=(\mathrm{FAIL},\mathrm{WEAK},\mathrm{FAIL},
 \mathrm{FAIL},\mathrm{FAIL}),
\]
overall \texttt{ROUTE\_A\_REJECTED}, Route B false, with scope
\texttt{NO\_BAD\_EULER\_OR\_ROOT\_NUMBER}.

\paragraph{Ownership and derivation ledger.}
The equality of the minimal ultimate matrix period with $\gamma$ is classical
and owned by the first source; the ultimate CSR representation is owned by the
second.  The package-level increment is their convention-consistent synthesis
with the one-equality transient proof, exact divisor strata, phase-indexed
ultimate spans, the fixed-support sharp family, and the reducible firewall.
The census validates these formulas on declared finite inputs but supplies no
priority or completeness claim about the wider literature.

\ifcase\CRevisionRound
\paragraph{Revision-round focus.} Round 0 freezes the irreducible rational
max-plus convention, exact critical cyclicity, first-equality transient and
the fixed-support family with every positive transient.
\or
\paragraph{Revision-round focus.} Round 1 adds the source-locked CSR phases,
raw and projective attraction strata, the eigencone, and the periodic family
of ultimate column spans.
\or
\paragraph{Revision-round focus.} Round 2 separates classical ownership from
finite regression, closes the reducible multirate boundary, and locks the
semantic mutation audit and five-gate Route-A stop.
\fi

\paragraph{Limitations and nonclaims.}
We claim no priority for cyclicity, attraction-cone, ultimate-span, or CSR
theory; no uniform transient; and no equality of every vector period with
$\gamma$.  We assert no arithmetic local data, target divisor,
Hilbert--P\'olya operator, global novelty, external review, or acceptance score.

\section*{Declarations}
\paragraph{Data and code availability.} Package-local exact evidence and
deterministic code accompany this manuscript.
\paragraph{Ethics.} No human, animal, clinical, personal, or sensitive data
are used; approval is not applicable.
\paragraph{Author contributions (CRediT).} This anonymous certificate records
Conceptualization, Formal analysis, Software, Validation, Writing---original
draft, and Writing---review and editing. AI systems are not authors.
\paragraph{Funding.} No external funding is reported.
\paragraph{Conflicts of interest.} None known.
\paragraph{AI-use disclosure.} An AI coding assistant supported drafting and
exact-code development; it was not an external reviewer or independent error
process.

\begin{thebibliography}{2}
\bibitem{Sergeev09} S. Sergeev,
``Max algebraic powers of irreducible matrices in the periodic regime: An
application of cyclic classes,'' \emph{Linear Algebra Appl.} 431(8) (2009),
1325--1339. DOI: 10.1016/j.laa.2009.04.027.

\bibitem{SS12} S. Sergeev and H. Schneider,
``CSR expansions of matrix powers in max algebra,''
\emph{Trans. Amer. Math. Soc.} 364(11) (2012), 5969--5994.
DOI: 10.1090/S0002-9947-2012-05605-4.
\end{thebibliography}
\end{document}
